% Acknowledgements moved into \Dedication in metadata.tex (printed on the
% dedication page, right after the signed declaration in title.tex).

\chapter*{Introduction}
\addcontentsline{toc}{chapter}{Introduction}
\label{chap:introduction}

Every project, from building a bridge to releasing a piece of software,
ultimately comes down to the same question: in what order should the work
get done? \changed{The answer can rarely be chosen freely.} Activities depend on one another, since
the foundation has to be poured before the walls go up, and the people,
machines and budget available to do the work are finite. A project
\emph{schedule} is a concrete answer to this question: an assignment of a
start time, and sometimes a choice of \emph{how} an activity is to be
carried out, to every piece of work in the project, consistent with both
kinds of constraint. Getting it right matters in a very direct, practical
sense, since a better schedule finishes the same project sooner, or finishes
it with the resources actually on hand, without changing what has to be
built or who is available to build it.

What makes this question hard is not any one constraint in isolation, but
the way they interact. Adding more activities, more dependencies, or more
ways of carrying out each activity does not make the problem grow
gracefully; it makes the space of possible schedules grow combinatorially,
to the point where even a modern computer cannot examine every option for
all but the smallest projects. This is true even in the simplest version of
the problem, where every activity has only one way of being done; allowing
multiple \emph{modes}, for instance finishing a task faster by paying for
more labour, or slower with less, only deepens the difficulty, since now a
schedule must also choose, for every activity, which mode to use. This is
the \emph{multi-mode resource-constrained project scheduling problem}
(MRCPSP), the subject of this thesis, and it is provably among the hardest
problems in combinatorial optimisation: no known algorithm can guarantee an
optimal schedule for realistically sized instances within a practical
amount of time.

This is precisely the kind of problem that \emph{evolutionary algorithms}
were designed for. Rather than searching for a guaranteed-optimal answer, an
evolutionary algorithm maintains and improves a population of candidate
solutions over many generations, borrowing the mechanics of natural
selection, namely variation, competition, and survival of the fitter, to
explore an otherwise unmanageable search space. They do not need the
problem to have any convenient mathematical structure, which makes them
well suited to a problem like the MRCPSP, where the constraints that make
exact methods struggle are exactly the constraints an evolutionary search
can simply evaluate, rather than solve, for every candidate it considers.
This thesis pursues one specific evolutionary approach in this family: a
\emph{genetic-programming hyper-heuristic}, which does not evolve a
schedule directly but instead evolves the \emph{rule} that builds one,
decision by decision, as the project unfolds.

The chapters that follow develop this idea. Chapter~\ref{chap:theoretical-background}
introduces the classical theory of project scheduling and the
single-mode resource-constrained problem it builds on.
Chapter~\ref{chap:overview} extends this to the multi-mode case and surveys
the evolutionary and exact methods applied to it, including
the baseline this thesis builds on directly. Chapter~\ref{chap:dataset}
describes the benchmark datasets and the metrics used to judge a schedule's
quality. Chapter~\ref{chap:approach}
presents the proposed algorithm, and Chapter~\ref{chap:experiments} reports
how it performs against the baseline. \nameref{chap:conclusion} closes the thesis with what was learned
and what remains open.
